\section*{الفصل \ref{ch:b2:precambrian-earth} --- الأرض فيما قبل الكمبري}

\begin{solution}{exo:b2:precambrian-earth:1}
(أ) \omterm{def:b2:precambrian-earth:eons}{الدهر الجهنمي} $4567 - 4031 = \qty{536}{Myr}$؛ \omterm{def:b2:precambrian-earth:eons}{والدهر السحيق} $4031 - 2500 = \qty{1531}{Myr}$؛
\omterm{def:b2:precambrian-earth:eons}{ودهر الحياة الطليعية} $2500 - 538.8 = \qty{1961}{Myr}$. (ب) $(4567 - 538.8)/4567 = 0.882$: أي نحو سبعة
أثمان. (ج) $4567 - 4404 = \qty{163}{Myr}$ للزركون؛ و$4567 - 4030 = \qty{537}{Myr}$
لأقدم الصخور.
\end{solution}

\begin{solution}{exo:b2:precambrian-earth:2}
(أ) العوامل $e^{\lambda t}$ مع $\lambda = \ln2/T$ و$t = \qty{3}{Gyr}$ هي 1.593
(\ce{^{238}U})، و19.18 (\ce{^{235}U})، و1.160 (\ce{^{232}Th}) و5.292 (\ce{^{40}K}): فكانت
الإسهامات 3.053 و1.573 و2.424 و\qty{4.329}{pW}، أي \qty{11.38}{pW} في المجموع، مقابل
\qty{4.906}{pW} اليوم. (ب) $11.38/4.906 = 2.32$. (ج) الثوريوم-232 يتصدّر اليوم؛ وكان البوتاسيوم-40 يتصدّر
حينئذ.
\end{solution}

\begin{solution}{exo:b2:precambrian-earth:3}
(أ) $1616 + 18.0\times(28 - 32.2) = \qty{1540}{\celsius}$. (ب) $32.2 + (1500 - 1616)/18.0 =
\qty{25.8}{\percent}$. (ج) تحليل الصخر الكلّي ليس تركيب سائل: فمسيل
\omterm{def:b2:precambrian-earth:greenstone}{الكوماتيت} قد يحوي بلّورات أوليفين متراكمة من السائل، تضيف أكسيد المغنيسيوم. وقيمته
\qty{38}{\percent} تتجاوز \qty{32}{\percent} في أغنى السوائل الكوماتيتية، وتتجاوز
التجارب التي يستند إليها الخطّ.
\end{solution}

\begin{solution}{exo:b2:precambrian-earth:4}
أقدم الصخور، \qty{4.03}{Ga}؛ أولى الستروماتوليتات، \qty{3.43}{Ga}؛ نهاية قباب بيلبارا،
\qty{3.1}{Ga}؛ \omterm{def:b2:precambrian-earth:goe}{حدث الأكسدة العظيم}، من نحو 2.4 إلى \qty{2.1}{Ga}؛ آخر الكوماتيتات الشائعة،
\qty{1.9}{Ga}؛ الارتفاع الثاني للأكسجين، نحو \qty{0.8}{Ga}؛ نهاية المارينوي،
\qty{0.635}{Ga}؛ أولى أحافير الإدياكاري، نحو \qty{0.57}{Ga}.
\end{solution}

\begin{solution}{exo:b2:precambrian-earth:5}
(أ) $\tfrac14\times1361\times0.706 = \qty{240.2}{W.m^{-2}}$ ممتصّة و$\sigma\times288^4 =
\qty{390.1}{W.m^{-2}}$، فيكون $\varepsilon = 240.2/390.1 = 0.616$. (ب) $\tfrac14\times1361\times0.4
= \qty{136.1}{W.m^{-2}}$، و$T^4 = 136.1/(0.616\times\num{5.670e-8})$، و$T = \qty{249.8}{K}$.
(ج) عند \qty{275.5}{K} يكون الألبيدو $0.6 - 0.306\times12.5/20 = 0.4088$: فيكون $R = 340.25\times0.5912 =
\qty{201.2}{W.m^{-2}}$ و$E = 0.616\,\sigma\times275.5^4 = \qty{201.2}{W.m^{-2}}$. (د) $R' =
340.25\times0.306/20 = \qty{5.21}{W.m^{-2}.K^{-1}}$، و$E' = 4\varepsilon\sigma T^3 =
\qty{2.92}{W.m^{-2}.K^{-1}}$: $R' > E'$، فالتوازن غير مستقرّ
(\cref{prop:b2:precambrian-earth:ebm}).
\end{solution}

\begin{solution}{exo:b2:precambrian-earth:6}
(أ) $1/(1 + 0.4\times2.4/4.567) = 0.826$؛ و$1/(1 + 0.4\times0.65/4.567) = 0.946$. (ب) تحتاج الحالة
الدافئة إلى 0.932 على الأقلّ من ضوء شمس اليوم: فهي مستحيلة قبل \qty{2.4}{Ga}، وممكنة
قبل \qty{0.65}{Ga}. (ج) $0.826/0.946 = 0.873$: أي أخفت بنسبة \qty{13}{\percent}، قريبًا من
\qty{10}{\percent} المذكورة، التي تأتي من نموذج شمسي مختلف.
\end{solution}

\begin{solution}{exo:b2:precambrian-earth:7}
(أ) الستيرتي $716.9 - 658.9 = \qty{58.0}{Myr}$؛ والمارينوي $644.5 - 635.4 = \qty{9.1}{Myr}$؛
والفترة بين الجليدية $658.9 - 644.5 = \qty{14.4}{Myr}$. (ب) $(58.0 + 9.1)/85 = 0.79$: أي نحو أربعة أخماس
العصر. (ج) بدايته محصورة فقط بين 649.9 و\qty{639.0}{Ma}: فبحسب هذين
الحدّين دام المارينوي بين نحو 3 و\qty{15}{Myr}.
\end{solution}

\begin{solution}{exo:b2:precambrian-earth:8}
(أ) $\tau = \frac{0.05}{0.95}\times\frac{44.01}{28.97}\times\frac{\num{5.1e18}}{\num{0.26e12}} =
0.05263\times1.519\times\num{1.962e7} = \qty{1.57}{Myr}$. (ب) $0.05/\num{278e-6} = 180$ ضعفًا.
(ج) ضعف المدّة، أي \qty{3.14}{Myr}.
\end{solution}

\begin{solution}{exo:b2:precambrian-earth:9}
(أ) سقطت الحجارة في طين طريّ فثنت صفائحه نحو الأسفل: إنها حجارة ساقطة، أطلقها
الجليد الطافي. (ب) الجليد الذي يبلغ البحر على بعد \ang{10} من خطّ الاستواء هو الأثر المتوقّع
للكرة الثلجية، ومن الصعب جدًا الحصول عليه إن لم تتجمّد إلا المناطق القطبية. (ج) كربونات غطائية:
تسجّل النهاية المفاجئة للتجلّد، مع تجوية شديدة وغمر مياه الذوبان
للأرصفة.
\end{solution}

\begin{solution}{exo:b2:precambrian-earth:10}
(أ) $0.94 > 0.932$، أي نقطة التحوّل الدافئة: فالتوازن الدافئ، $288\times0.94^{1/4} =
\qty{283.6}{K}$، يقع فوق \qty{283}{K} بقليل؛ و$0.94 < 1.227$، فالكرة الثلجية موجودة أيضًا.
(ب) $249.8\times0.94^{1/4} = \qty{246.0}{K}$. (ج) $\varepsilon = 0.94\times1361\times0.4/(4\sigma
\times263^4) = 511.7/1085.1 = 0.472$. (د) $\sigma\times263^4 = \qty{271.3}{W.m^{-2}}$، و$(0.616
- 0.472)\times271.3 = \qty{39}{W.m^{-2}}$.
\end{solution}

\begin{solution}{exo:b2:precambrian-earth:11}
(أ) $S_{\text{m}}/S_{\text{f}} = (T_1/T_2)^4\,(1 - a_{\text{w}})/(1 - a_{\text{i}})$: فالانبعاثية
تختزل. $(263/283)^4\times0.706/0.4 = 0.7459\times1.765 = 1.317$. (ب) تباين أصغر
في الألبيدو بين الجليد والماء المفتوح، كجليد أدكن مغبرّ، أو منحدر أعرض في
درجة الحرارة. (ج) تعكس الكرة الثلجية قدرًا كبيرًا من ضوء الشمس فتبقى باردة تحت أي ضوء شمس أو
دفيئة كانت الأرض الدافئة ستتجمّد عندهما؛ ولإذابتها يجب أن يتجاوز التأثير تأثير
التجمّد بكامل عرض الحلقة.
\end{solution}

\begin{solution}{exo:b2:precambrian-earth:12}
قرب $T^*$، $C\dot\delta = -(E' - R')\delta$، فيكون $\delta \propto e^{-t/\tau}$ مع $\tau =
C/(E' - R')$. وعند التوازن الدافئ $R' = 0$ و$E' = 4\times0.616\times\num{5.670e-8}\times
288^3 = \qty{3.34}{W.m^{-2}.K^{-1}}$؛ و$C = 1025\times3992\times50 = \qty{2.05e8}{J.m^{-2}.K^{-1}}$،
فيكون $\tau = \qty{6.1e7}{s}$، أي نحو \qty{1.9}{yr}. وفي هذا النموذج يبقى عند التوازن الدافئ $R' =
0$ حتى يبلغ \qty{283}{K}، فيبقى $\tau$ قرب السنتين حتى نقطة التحوّل، ثم
تختفي الحالة. ومع ألبيدو يتغيّر تغيّرًا سلسًا، ينمو $R'$ كلّما اقترب التوازن من
الانطواء، حيث $E' = R'$: فينمو $\tau$ بلا حدّ، وهو تباطؤ ينذر بنقطة
التحوّل.
\end{solution}

\begin{solution}{pb:b2:precambrian-earth:1}
\textbf{1.} $\num{0.0203e-6}\times0.9929\times\num{95.13e-6} = \qty{1.92}{pW.kg^{-1}}$
(\ce{^{238}U})؛ و$\num{0.0203e-6}\times0.00711\times\num{568.47e-6} = \qty{0.0820}{pW.kg^{-1}}$
(\ce{^{235}U})؛ و$\num{0.0795e-6}\times\num{26.28e-6} = \qty{2.09}{pW.kg^{-1}}$ (\ce{^{232}Th})؛
و$\num{240e-6}\times\num{1.20e-4}\times\num{28.47e-6} = \qty{0.820}{pW.kg^{-1}}$ (\ce{^{40}K}).
\textbf{2.} \qty{4.91}{pW.kg^{-1}}.
\textbf{3.} $\ln2/T$: 0.155 و0.985 و0.0495 و\qty{0.555}{Gyr^{-1}}.
\textbf{4.} العوامل 1.52 و14.3 و1.14 و4.48: فالإسهامات 2.91 و1.17 و2.39 و\qty{3.67}{pW.kg^{-1}}.
\textbf{5.} \qty{10.1}{pW.kg^{-1}}، أي 2.07 ضعف قيمة اليوم.
\textbf{6.} البوتاسيوم-40. فعمر نصفه، \qty{1.25}{Gyr}، قصير: فقد مضى منذ ذلك الحين أكثر من عمري نصف،
وتفكّك أكثر من ثلاثة أرباع \ce{^{40}K} الذي كان موجودًا حينئذ.
\textbf{7.} $1/(1 + 0.4\times2.7/4.567) = 0.809$؛ و$0.809\times1361 = \qty{1101}{W.m^{-2}}$.
\textbf{8.} $288\times0.809^{1/4} = \qty{273}{K}$، أي دون \qty{283}{K}: فليست حالة دافئة صالحة؛
وكانت الأرض ستتجمّد.
\textbf{9.} $(1101\times0.4/(4\times0.616\,\sigma))^{1/4} = \qty{237}{K}$.
\textbf{10.} $0.616\times0.809 = 0.498$.
\textbf{11.} من $0.004/\num{278e-6} = 14$ إلى $0.03/\num{278e-6} = 108$ ضعفًا.
\textbf{12.} المنظّم الحراري لتجوية السيليكات: تتباطأ التجوية في المناخ البارد، فتدع
ثاني أكسيد الكربون البركاني يتراكم، وتتسارع في المناخ الدافئ.
\textbf{13.} 0.826 من ضوء اليوم، أي دون 0.932: فلا حالة دافئة؛ وكانت الأرض بهواء اليوم ستكون
كرة ثلجية.
\textbf{14.} $\frac{0.1}{0.9}\times1.519\times\num{1.962e7} = \qty{3.31}{Myr}$.
\textbf{15.} $\frac{0.2}{0.8}\times1.519\times\num{1.962e7} = \qty{7.45}{Myr}$.
\textbf{16.} \omterm{def:b2:precambrian-earth:snowball}{الكربونات الغطائية}، وهي دولوستون مصفّح يعلو فجأةً الرواسب الجليدية،
ترسّب من بحر غذّته تجوية شديدة في المرحلة اللاحقة الحارّة وغمرته
مياه الذوبان.
\textbf{17.} لا يبقى الحديد (II) ذائبًا إلا في ماء بلا أكسجين: فيتراكم الحديد في
المحيط الخالي من الأكسجين تحت الجليد ويترسّب على هيئة حديد (III) حيث يصل الأكسجين، مثلًا مع
مياه الذوبان.
\textbf{18.} $(1616 - 1450)/(32.2 - 23) = \qty{18.0}{\celsius}$ لكل واحد في المئة.
\textbf{19.} $1616 + 18.0\times(30 - 32.2) = \qty{1576}{\celsius}$.
\textbf{20.} \qty{1616}{\celsius}، أي النقطة المقيسة.
\textbf{21.} $1616 - 18.0\times20.2 = \qty{1252}{\celsius}$.
\textbf{22.} $1616 - 20\times20.2 = \qty{1212}{\celsius}$: أخفض بمقدار \qty{40}{\celsius}، أي عُشر
الفرق بين نوعي الحمم؛ فالخلاصة لا تتوقّف على الاختيار.
\textbf{23.} جاءت الكوماتيتات من أعمدة، أسخن من الوشاح المحيط، ومن درجات انصهار كبيرة
في العمق؛ وحرارة الحمم ليست درجة الحرارة الكامنة لمصدرها. أمّا
الوشاح ككلّ فكان أدفأ بمقدار 200 إلى \qty{300}{K}.
\textbf{24.} كانت أسخن حمم كوماتيتية أسخن من بازلت أوّلي لظهرات اليوم بنحو
\textbf{\qty{364}{\celsius}}.
\end{solution}
